\section{生成函数}

许多数理物理中的函数集可用生成函数定义。这类函数包括但不限于
我们在罗德里格斯公式讨论中涉及的正交多项式 $y_n$。目前，我们不
对函数的来源做假设。

若 $f_n(x)$ 是以整数 $n$ 为指标的函数集，则这些 $f_n(x)$ 可描述为
某辅助变量 $t$ 的幂级数展开中某函数 $g(x, t)$ 的系数，这个
$g(x, t)$ 就称为生成函数：

$$
g(x, t)=\sum_n c_n f_n(x) t^n
$$

这里 $n$ 的取值可为半无限（$n \geq 0$，即泰勒级数），也可为
$-\infty$ 到 $+\infty$（即劳伦级数）。额外的系数 $c_n$ 用于调整
归一化。对同一组 $f_n$，不同 $c_n$ 会导致不同的生成函数 $g(x, t)$。

应用留数定理，生成函数的展开与 $f_n$ 的复变积分表示相关：

$$
c_n f_n(x)=\frac{1}{2 \pi i} \oint \frac{g(x, t)}{t^{n+1}} d t
$$

其中积分路径围绕 $t=0$，但不包围被积函数（关于 $t$）的其他奇点。

生成函数既可用来定义 $f_n(x)$，也可作为已知 $f_n$ 的封装（如
$f_n$ 是 Sturm-Liouville 常微分方程的多项式解）。我们稍后讨论如何
为已知 $f_n$ 构造生成函数，目前只关注生成函数的用法。

显然，通过显式展开生成函数，可提取函数集的各成员。更重要的是，
生成函数在推导 $f_n$ 集合成员间关系时非常有用。例如，

$$
\frac{\partial g(x, t)}{\partial t}=\sum_n n c_n f_n(x) t^{n-1}=
\sum_n(n+1) c_{n+1} f_{n+1}(x) t^n,
$$

若能将 $g$ 与 $\partial g / \partial t$ 联系，就能得到 $f_n$ 与
$f_{n+1}$ 的关系。对 $g(x, t)$ 关于 $x$ 求导，还可推导 $f_n(x)$
\begin{equation}\label{eq:hermite-gen}
e^{-t^2+2 t x}=\sum_{n=0}^{\infty} H_n(x) \frac{t^n}{n!}
\end{equation}
$$
e^{-t^2+2 t x}=\sum_{n=0}^{\infty} H_n(x) \frac{t^n}{n!}
$$

为推导 $H_n$ 的递推公式，计算

$$
\begin{equation}\label{eq:hermite-dt}
\frac{\partial}{\partial t} e^{-t^2+2 t x}=(2 x-2 t) e^{-t^2+2 t x}=\sum_{n=0}^{\infty} n H_n(x) \frac{t^{n-1}}{n!} .
\end{equation}
将指数函数在等式\eqref{eq:hermite-dt}中间项展开（暂省 $H_n$ 的自变量），有

$$
\sum_{n=0}^{\infty} 2 x H_n \frac{t^n}{n!}-\sum_{n=0}^{\infty} 2 H_n
\frac{t^{n+1}}{n!}=\sum_{n=0}^{\infty} n H_n \frac{t^{n-1}}{n!}
$$

从每项中提取 $t^n$ 的系数，得（对每个 $n$）

$$
\frac{2 x H_n}{n!}-\frac{2 H_{n-1}}{(n-1)!}=\frac{(n+1) H_{n+1}}{(n+1)!},
$$

化简为

$$
2 x H_n(x)-2 n H_{n-1}(x)=H_{n+1}(x) .
$$

\begin{equation}\label{eq:hermite-recur}
2 x H_n(x)-2 n H_{n-1}(x)=H_{n+1}(x) .
\end{equation}

$$
\frac{\partial}{\partial x} e^{-t^2+2 t x}=2 t e^{-t^2+2 t x}=
\sum_{n=0}^{\infty} H_n^{\prime}(x) \frac{t^n}{n!} .
$$

将等式\eqref{eq:hermite-gen}代入该式中间项，得

$$
\sum_{n=0}^{\infty} 2 H_n(x) \frac{t^{n+1}}{n!}=
\sum_{n=0}^{\infty} H_n^{\prime}(x) \frac{t^n}{n!},
$$

直接得到

$$
2 n H_{n-1}(x)=H_n^{\prime}(x) .
$$


\subsection{寻找生成函数}

为了使生成函数不再显得神秘，下面我们来讨论它们是如何被构造出来的。
对于或多或少任意的一类函数而言，这一问题一直是数学研究中的一个活跃课题。
在过去的一个世纪中，Rainville、Weisner、Truesdell 等人发展出了多种不同的方法。
有关这方面的工作，可参见“补充阅读”中 McBride 和 Talman 的著作。
\begin{equation}\label{eq:gen-schlaefli}
g(x, t)=\frac{1}{w(x)} \sum_{n=0}^{\infty} c_n t^n \frac{n!}{2 \pi i} \oint_C \frac{w(z)[p(z)]^n}{(z-x)^{n+1}} d z
\end{equation}

$$
g(x, t)=\frac{1}{w(x)} \sum_{n=0}^{\infty} c_n t^n \frac{n!}{2 \pi i} \oint_C \frac{w(z)[p(z)]^n}{(z-x)^{n+1}} d z
$$

记住 $C$ 包含 $x$，且 $w p^n$ 在围道所圈定的区域内必须解析。

原则上，方程 \eqref{eq:gen-schlaefli} 可以被求值以得到 $g(x,t)$，例如可以选择适当的 $C$ 使得求和可移入 $z$ 积分内，然后（在指定 $c_n$ 之后）先求和再评价围道积分。实际上，能否顺利完成此步骤取决于具体问题及 $c_n$ 的选择。下面给出一个该过程的示例。

我们用上述形式化过程来得到勒让德多项式的生成函数。勒让德微分方程具有式 (12.1) 中讨论的形式，

$$
\left(1-x^2\right) y^{\prime \prime}-2 x y^{\prime}+\lambda y=0,
$$

因此

$$
p(x)=1-x^2, \quad q(x)=-2 x,
$$

并且该方程按所写形式是自共轭的，所以 $w(x)=1$。根据基于 Schlaefli 积分的生成函数公式（方程 (12.25)），我们选择 $c_n=(-1)^n / 2^n n!$，从而得到

$$
g(x, t)=\sum_{n=0}^{\infty}\left(\frac{(-1)^n t^n}{2^n n!}\right) \frac{n!}{2 \pi i} \oint_C \frac{\left(1-z^2\right)^n}{(z-x)^{n+1}} d z
$$

将求和與積分互換（下文将对此进行说明），依赖于 $n$ 的因子构成一个几何级数，可以合并求和：

\begin{equation}\label{eq:geo-sum}
\begin{aligned}
\sum_{n=0}^{\infty}\left(\frac{\left(z^2-1\right) t}{2(z-x)}\right)^n \frac{1}{z-x} & =\frac{1}{z-x-\frac{1}{2}\left(z^2-1\right) t} \\
& =-\frac{2}{t}\left[z^2-\frac{2 z}{t}+\frac{2 x-t}{t}\right]^{-1}
\end{aligned}
\end{equation}


将该结果代回 $g(x,t)$ 的表达式，得到

$$
\begin{aligned}
g(x, t) & =-\frac{2}{t} \frac{1}{2 \pi i} \oint_C\left[z^2-\frac{2 z}{t}+\frac{2 x-t}{t}\right]^{-1} d z \\
& =-\frac{2}{t} \frac{1}{2 \pi i} \oint_C \frac{d z}{\left(z-z_1\right)\left(z-z_2\right)}
\end{aligned}
$$

其中 $z_1$ 和 $z_2$ 是第一行括号内二次形式的根：

$$
z_1=\frac{1}{t}-\frac{\sqrt{1-2 x t+t^2}}{t}, \quad z_2=\frac{1}{t}+\frac{\sqrt{1-2 x t+t^2}}{t} .
$$

为了使方程 \eqref{eq:geo-sum} 有效，必须证明将求和与积分互换是合法的，这只有在对所有在围道 $C$ 上使用的点（即围道上的任意点）求和对 $z$ 是一致收敛时才成立。便于分析收敛性的一个情形是取小的 $t$ 和 $x$ 并选取 $|z|=1$ 的围道。一旦得到最终公式，可通过解析延拓来扩大其有效域。

在所假定的围道和小的 $x$ 下，存在一段 $|t|\ll 1$ 的范围使得

$$
\left|\frac{\left(z^2-1\right) t}{2(z-x)}\right|<1,
$$

从而保证几何级数的收敛。现在回到对方程 \eqref{eq:geo-sum} 中的围道积分的评价。被积函数在 $z=z_1$ 和 $z=z_2$ 处有两个极点。对小的 $x$ 与小的 $|t|$，$z_2$ 约为 $2/t$，位于围道之外，而 $z_1$ 则靠近 $z$ 平面的原点。因此，只有位于 $z=z_1$ 处的留数对围道积分有贡献，该积分的值为

$$
g(x, t)=-\frac{2}{t} \frac{1}{z_1-z_2} .
$$

由于

$$
z_1-z_2=-\frac{2}{t} \sqrt{1-2 x t+t^2}
$$

我们得到勒让德多项式的生成函数为

$$
g(x, t)=\frac{1}{\sqrt{1-2 x t+t^2}}
$$

对于由 Sturm-Liouville 问题产生并可由 Rodrigues 公式描述的一类多项式，我们可以更具体地说明。利用 Schlaefli 积分（方程 (12.18)），我们可以构造

$$
g(x, t)=\frac{1}{w(x)} \sum_{n=0}^{\infty} c_n t^n \frac{n!}{2 \pi i} \oint_C \frac{w(z)[p(z)]^n}{(z-x)^{n+1}} d z
$$

记住 $C$ 包含 $x$，且 $w p^n$ 在围道所圈定的区域内必须解析。

原则上，方程 (12.25) 可以被求值以得到 $g(x,t)$，例如可以选择适当的 $C$ 使得求和可移入 $z$ 积分内，然后（在指定 $c_n$ 之后）先求和再评价围道积分。实际上，能否顺利完成此步骤取决于具体问题及 $c_n$ 的选择。下面给出一个该过程的示例。

我们用上述形式化过程来得到勒让德多项式的生成函数。勒让德微分方程具有式 (12.1) 中讨论的形式，

$$
\left(1-x^2\right) y^{\prime \prime}-2 x y^{\prime}+\lambda y=0,
$$

因此

$$
p(x)=1-x^2, \quad q(x)=-2 x,
$$

并且该方程按所写形式是自共轭的，所以 $w(x)=1$。根据基于 Schlaefli 积分的生成函数公式（方程 \eqref{eq:gen-schlaefli}），我们选择 $c_n=(-1)^n / 2^n n!$，从而得到

$$
g(x, t)=\sum_{n=0}^{\infty}\left(\frac{(-1)^n t^n}{2^n n!}\right) \frac{n!}{2 \pi i} \oint_C \frac{\left(1-z^2\right)^n}{(z-x)^{n+1}} d z
$$

将求和与积分互换（下文将对此进行说明），依赖于 $n$ 的因子构成一个几何级数，可以合并求和：

$$
\begin{aligned}
\sum_{n=0}^{\infty}\left(\frac{\left(z^2-1\right) t}{2(z-x)}\right)^n \frac{1}{z-x} & =\frac{1}{z-x-\frac{1}{2}\left(z^2-1\right) t} \\
& =-\frac{2}{t}\left[z^2-\frac{2 z}{t}+\frac{2 x-t}{t}\right]^{-1}
\end{aligned}
$$

将该结果代回 $g(x,t)$ 的表达式，得到

$$
\begin{aligned}
g(x, t) & =-\frac{2}{t} \frac{1}{2 \pi i} \oint_C\left[z^2-\frac{2 z}{t}+\frac{2 x-t}{t}\right]^{-1} d z \\
& =-\frac{2}{t} \frac{1}{2 \pi i} \oint_C \frac{d z}{\left(z-z_1\right)\left(z-z_2\right)}
\end{aligned}
$$

其中 $z_1$ 和 $z_2$ 是第一行括号内二次形式的根：

$$
z_1=\frac{1}{t}-\frac{\sqrt{1-2 x t+t^2}}{t}, \quad z_2=\frac{1}{t}+\frac{\sqrt{1-2 x t+t^2}}{t} .
$$

为了使方程 \eqref{eq:geo-sum} 有效，必须证明将求和与积分互换是合法的，这只有在对所有在围道 $C$ 上使用的点（即围道上的任意点）求和对 $z$ 是一致收敛时才成立。便于分析收敛性的一个情形是取小的 $t$ 和 $x$ 并选取 $|z|=1$ 的围道。一旦得到最终公式，可通过解析延拓来扩大其有效域。

在所假定的围道和小的 $x$ 下，存在一段 $|t|\ll 1$ 的范围使得

$$
\left|\frac{\left(z^2-1\right) t}{2(z-x)}\right|<1,
$$

从而保证几何级数的收敛。现在回到对方程 (12.26) 中的围道积分的评价。被积函数在 $z=z_1$ 和 $z=z_2$ 处有两个极点。对小的 $x$ 与小的 $|t|$，$z_2$ 约为 $2/t$，位于围道之外，而 $z_1$ 则靠近 $z$ 平面的原点。因此，只有位于 $z=z_1$ 处的留数对围道积分有贡献，该积分的值为

$$
g(x, t)=-\frac{2}{t} \frac{1}{z_1-z_2} .
$$

由于

$$
z_1-z_2=-\frac{2}{t} \sqrt{1-2 x t+t^2}
$$

我们得到勒让德多项式的生成函数为

$$
g(x, t)=\frac{1}{\sqrt{1-2 x t+t^2}}
$$
